% !TeX root = main_long.tex
\section{Introduction}

Healthcare research and personalized medicine depend on integrating clinical observations, molecular data, and biomedical knowledge.
These resources are distributed across institutions and databases, where differences in data structures, identifiers, and terminology complicate their combined use~\cite{gaudetblavignac2021sphn,putman2024monarch}.
Semantic Web approaches provide a basis for addressing this heterogeneity through shared identifiers, explicit relationships and reusable vocabularies~\cite{smith2007obo,belleau2008bio2rdf}.
Leveraging these strategies has enabled harmonizing clinical data across institutions while integrating knowledge about genes, diseases, and phenotypes to support biomedical investigations~\cite{gaudetblavignac2021sphn,putman2024monarch}.
Meaningful integration of these distributed resources nevertheless depends on retaining the context in which source observations were recorded, not only on adopting common identifiers or terminology~\cite{gaudetblavignac2021sphn}.
For genomic data, this motivates a reusable semantic representation that preserves source observations and their interpretation context while enabling connections to clinical information and biomedical knowledge.

The Variant Call Format (VCF) is widely used to exchange genomic variant data~\cite{danecek2011vcf}.
A VCF file can be understood as a structured table in which each data row describes a genomic location and the sequence observed there relative to a reference genome.
For each row, a record contains the reference allele, one or more alternative alleles, and optional annotations describing properties such as quality, filtering, or predicted biological effects.
When individual samples are included, additional columns record sample-level observations such as the genotype or sequencing depth.
A header at the beginning of the file defines many of these fields and specifies how their values should be interpreted~\cite{vcf45}.

VCF format has been evolving since its first public release in 2011.
At the time of writing, the specifications are maintained by the Global Alliance for Genomics and Health (GA4GH) \cite{ga4gh2016federated} with versions 4.1--4.5 hosted on a GA4GH-maintained github repository~\cite{ga4gh_hts_specs}.
The most current specification is VCF v4.5, released February 25, 2026~\cite{vcf45}.
VCF 4.1--4.5 share the basic structures described above but differ in supported declarations and constraints~\cite{htsspecs,vcf45}. 
For example, later versions introduce new ways to describe variants and their associated metadata, while also refining how existing fields should be interpreted. 
% TODO: need transition sentence
These relationships are important because VCF often encodes meaning indirectly through the structure of the file.
For example, the reference allele is assigned index 0 and alternative alleles are numbered in the order in which they appear.
A genotype written as \term{0/2} therefore does not itself identify two DNA sequences: it indicates that the sample carries the reference allele and the second listed alternative allele.
Similarly, the meaning and expected number of values for an annotation may depend on a definition given in the file header.
This simple example illustrates a broader characteristic of VCF: the format achieves compactness by relying on relationships that are understood within the file rather than stated explicitly with each value. 
More complex \emph{cases} depend on several such relationships at once. 
When VCF observations are integrated with information outside the file, these implicit relationships must remain available if the observations are to be integrated meaningfully with other genomic, clinical, and biomedical data.
The challenge is therefore not that VCF lacks contextual information, but that much of this context is encoded through the internal structure of the file rather than exposed as explicit, independently addressable relationships.

A semantic representation offers a way to make these dependencies explicit while preserving their connection to the original observations.
Doing so supports semantic interoperability and reuse, central concerns of the Findable, Accessible, Interoperable and Reusable (FAIR) principles~\cite{wilkinson2016fair}.
The wider FAIRness of a dataset also depends on its identifiers, metadata, access conditions and provenance, rather than on the choice of file format alone~\cite{bahim2020fair}.
The Resource Description Framework (RDF) provides a graph-based model for representing resources and their relationships, while Linked Data principles support the use of shared identifiers and links across datasets~\cite{decker2000semantic,bizer2009linked}.
Together, these provide a basis for exposing VCF observations and their interpretation context alongside connections to external biomedical resources, which is necessary for any downstream clinical or research applications \cite{lubin2017variantfile}.

\subsection{Related Work}
Related work includes reusable semantic models and VCF--RDF conversion or query tools, with some contributions combining both roles (Table~\ref{tab:vcf-rdf-related-work}).
The Genomic Feature and Variation Ontology (GFVO), Genome Variation Ontology (GVO), and HERO-Genomics model genomic or clinical concepts; tools such as VCF2RDF, TogoVar, BioInterchange, AgroLD and bio-vcf~\cite{garrison2022vcflib} expose VCF-derived data for integration.
The GA4GH Variation Representation Specification (VRS) complements these approaches with computable variation descriptions and federated identification, using JSON Schema rather than RDF or OWL; it does not preserve source-file structure~\cite{wagner2021vrs}.

\begin{table}[htbp]
\caption{VCF-related approaches and their roles: reusable models, RDF conversion/query tools, or combinations of both.}
\label{tab:vcf-rdf-related-work}
\centering
\small
\begin{tabularx}{\linewidth}{@{}>{\raggedright\arraybackslash}p{0.27\linewidth}>{\raggedright\arraybackslash}X@{}}
\toprule
Approach & Role and focus \\
\midrule
BioHackathon \textit{vcf2rdf}~\cite{katayama2019biohackathon}
& \textbf{Converter.} Early VCF--RDF conversion, using FALDO for positions and tool-specific terms for other VCF concepts.
\\
GFVO~\cite{baran2015gfvo}
& \textbf{OWL ontology.} Harmonizes genomic features and variation across VCF, GVF, GFF3 and GTF.
\\
FALDO / BioInterchange~\cite{bolleman2016faldo,baran2015gfvo}
& \textbf{Ontology + converter.} BioInterchange reuses FALDO for genomic locations in VCF--RDF conversion; FALDO does not model VCF structure or fields.
\\
VCF2RDF~\cite{penha2017vcf2rdf}
& \textbf{Mapping/converter.} Isomorphic VCF--RDF mapping exposing file components independently with source context; not a general genomic ontology.
\\
TogoVar \textit{vcf2rdf}~\cite{mitsuhashi2022togovar}
& \textbf{Converter.} Simple RDF representation of population allele-frequency VCF data for integration with TogoVar's biomedical knowledge graphs.
\\
GVO~\cite{kawashima2023gvo}
& \textbf{OWL ontology.} Describes biological variation, particularly complex structural variants, rather than VCF structure and interpretation rules.
\\
VariantKG~\cite{prasanna2025variantkg}
& \textbf{Converter + application ontology.} Builds large VCF-derived graphs for variant analysis and graph machine learning, rather than VCF specification modelling.
\\
SPARQLing genomics~\cite{janssen2020sparqlinggenomics}
& \textbf{Converter + query platform.} Generates RDF from VCF and other genomic sources and serves it through a query interface; its terms are tool-specific rather than a model of the VCF specification.
\\
Semantic Beacons~\cite{bodrug2025semantic}
& \textbf{Integration mapping.} Maps GA4GH Beacon responses to RDF with FALDO, SO and GENO for federated genomic querying; not a VCF--RDF representation.
\\
HERO-Genomics~\cite{cazzaro2025hero}
& \textbf{Ontology.} Integrates VCF- and Beacon-derived information across clinical and research settings, with broader scope than VCF specification coverage.
\\
AgroLD VCF2RDF~\cite{larmande2025agrold}
& \textbf{Converter.} Incorporates VCF variation into AgroLD's plant-science knowledge graph; not a dedicated VCF ontology.
\\
VCF-RDFizer~\cite{crum2025semantifying}
& \textbf{Converter + precursor vocabulary.} Our earlier VCF--RDF workflow; VCF Core separates its model from conversion and narrows its scope to the VCF specification.
\\
\bottomrule
\end{tabularx}
\end{table}


These approaches distinguish source observations from biological interpretations.
VCF Core provides a shared, version-aware representation of source observations with links to existing genomic models, allowing converters and queries to reuse VCF interpretation rules.

\paragraph{The VCF Core Vocabulary.}
Here, we present and assess the \href{https://ecrum19.github.io/vcf-core-vocabulary/}{VCF Core Vocabulary} (\textit{vcfcv})~\cite{vcfc210}.
It combines source-context preservation, \textit{expanded} and \textit{condensed} sample representations, 123 curated alignments to existing models, and version-specific artifacts for VCF 4.1--4.5.
Four questions structure its assessment: 
\textbf{RQ1:} does VCF Core preserve the intended meaning of the VCF constructs it claims to support?
\textbf{RQ2:} do the version-specific artifacts and the two \emph{sample profiles} behave as documented?
\textbf{RQ3:} does the resulting representation support a reproducible integration task while retaining source context?
\textbf{RQ4:} how does VCF Core relate to other existing VCF semantic models?

Our results answer these questions as follows.
\textbf{RQ1:} VCF Core is show to (a) directly conform to roughly half of all extracted VCF specification \emph{requirements} from VCF 4.1--4.5 with gaps coming from untested \emph{requirements} rather than failures, and (b) a structured representation for every construct in a curated inventory of VCF 4.5. Thus, we assert that we cover a large portion of the \emph{requirements} in VCF specifications 4.1--4.5 and all the terms defined in our vocabulary can be mapped back to the original VCF specificaiton constructs.
\textbf{RQ2:} the version-specific registries reproduce declarations extracted from the specification texts; a second, separately implemented converter answers the same questions the same way on nearly all checks; the two \emph{sample profiles} agree on every question above the sample level and diverge, by design, only below it; and a file's logical content can be rebuilt almost entirely from the RDF alone, without falling back on any preserved copy of the original text.
\textbf{RQ3:} an end-to-end allele-dependent integration example returns the \emph{expected answers} with their source provenance.
\textbf{RQ4:} inspecting four published VCF-related models against nine retrieval criteria shows that VCF Core's distinguishing role is not broader modelling of genomic variation, but making the interpretation context a VCF file carries about itself directly addressable as graph structure rather than leaving each application to reconstruct it.

These assessments deliberately separate the existence of a representational mechanism from its implementation and from the correctness of the resulting data.
We then discuss how this foundation could support reusable conversion engines and integrated, medically-relevant genomic querying.
Throughout, \term{vcfc:} denotes \url{https://w3id.org/vcf-core/vocab#}.


% TODO: (maybe) Motivation idea --> complex variant info is interesting for medical and research uses but it requires the attached metadata and ... to be useful. Here we want to model those fields semantically to made semantic VCF data more useful for downstream applications.

